\renewcommand{\theHequation}{C.\arabic{equation}}\renewcommand{\theHfigure}{C.\arabic{figure}}\renewcommand{\theHtable}{C.\arabic{table}}
\section*{Приложение В. Дополнительные результаты и неопределённость}\label{app:c}
\renewcommand{\thefigure}{В.\arabic{figure}}\setcounter{figure}{0}
\renewcommand{\thetable}{В.\arabic{table}}\setcounter{table}{0}
\renewcommand{\theequation}{В.\arabic{equation}}\setcounter{equation}{0}
\subsection*{Вероятностные прогнозы и модель времени}
Для сравнения прогнозов вероятности использованы Brier score и логарифмическая потеря. На известных исходах с меткой $Z_{ij}$ и прогнозом $\widehat p_{ij}$ они равны
\begin{equation}\label{eq:scores}
\mathrm{BS}_j=\frac1{N_j}\sum_i(\widehat p_{ij}-Z_{ij})^2,\qquad
\mathrm{LL}_j=-\frac1{N_j}\sum_i\{Z_{ij}\log\widehat p_{ij}+(1-Z_{ij})\log(1-\widehat p_{ij})\}.
\end{equation}
Здесь $N_j$ --- число проверенных уравнений данного профиля. Для двух приближений Ремизова Brier score равен $0{,}116741$ и $0{,}046749$, для Монте-Карло --- $0{,}004030$. Низкое последнее значение при отсутствии положительных исходов не доказывает умения распознавать успешные случаи Монте-Карло. Дополнительные показатели и отдельный стохастический результат приведены ниже.


Графики распределений ошибки, соотношения точности и стоимости, времени и решений модели уже приведены в основной части (рис.~\ref{fig:ecdf}, \ref{fig:cost}, \ref{fig:times}, \ref{fig:choice}) и здесь не дублируются. Полноту картины дополняют показатели прогнозов в табл.~\ref{tab:forecasts}. Они получены на сохранённых исходах без калибровки по финальной выборке.
\begin{table}[!ht]\centering\caption{Качество замороженных прогнозов на известных финальных исходах}\label{tab:forecasts}
\small\begin{tabular}{lrrrr}\toprule Приближение&Задач&Brier&Logloss&Медиана $t/\widehat t$\\\midrule
Ремизов, $n=128$&60&0,116741&0,380690&0,973604\\
Ремизов, $n=256$&60&0,046749&0,196481&0,974990\\
Монте-Карло, $M=4096$&20&0,004030&0,064462&0,952419\\\bottomrule\end{tabular}\end{table}
Для времени также вычислялась средняя абсолютная логарифмическая ошибка
\begin{equation}
D_j=\frac1{N_j}\sum_{i=1}^{N_j}\left|\log\frac{t_{ij}}{\widehat t_{ij}}\right|.
\end{equation}
Она составила $0{,}026722$, $0{,}027978$ и $0{,}320285$ в порядке строк таблицы. Следовательно, при близкой медиане отношения времён Монте-Карло имеет значительно больший разброс ошибок прогноза. Одна новая реализация времени не позволяет отделить повторную изменчивость конкретного вызова от переноса модели на новое уравнение. Сохранённые диапазоны по остаткам семейной перекрёстной проверки описывают модельный разброс и не являются сертифицированными интервалами покрытия.

В обучающих данных 12 точных стохастических реализаций сосредоточены лишь в трёх семействах из 40; 180 реализаций не равнозначны 180 различным уравнениям. Это ограничивает перенос прогноза редких успехов.

Пять групп признаков имеют размеры $128+56+6+48+15=253$. В число аналитических ориентиров входит $\widetilde T=\max\{0,\log[2/(\delta10^{-8})]/\delta\}$ и масштабы $\sqrt{2\overline a\widetilde T}$, $B\widetilde T$, где $\delta=\lambda-C$. Это заранее вычислимые признаки, а не фактическое время интегрирования $T$, которое выбиралось с учётом правой части. Они не являются признаками точного решения.

Для полноты воспроизведения выбранной вероятностной модели приведены параметры \eqref{eq:calibration}. Калибраторы использовали только разрешённые обучающие семейства; финальные исходы в их подборе не участвовали.
\begin{table}[!ht]\centering\caption{Параметры сигмоидальной калибровки}\label{tab:calibration}
\begin{tabular}{lrr}\toprule Приближение&$\alpha_j$&$\beta_j$\\\midrule
Ремизов, $n=128$&0,6887102701&0,1279481655\\
Ремизов, $n=256$&0,3649140499&0,5706520816\\
Монте-Карло&0,3699116924&$-1{,}1887570760$\\\bottomrule\end{tabular}\end{table}
В ходе разработки также рассматривался прогноз нормированного логарифма ошибки с эмпирическим распределением остатков. Он сохранён как исследованный альтернативный вариант, но в окончательном правиле \eqref{eq:policy} используются именно сигмоидально откалиброванные вероятности и отдельные регрессии времени. Не заявляется, что выбранный алгоритм превосходит все возможные модели.

\subsection*{Парность сравнения и пределы статистических выводов}
Число индивидуальных потерь точных доставок относительно каждой постоянной стратегии равно нулю. Наблюдаемые условные доли потерь --- $0/53$ для $n=128$ и $0/57$ для $n=256$; их нельзя автоматически записывать как условный риск с одинаковым знаменателем 60. Внутренний критерий выбора политики при разработке использовал реальные доли успеха Монте-Карло, а не искусственно бинаризованные удачные реализации.

В описательном семейном повторном отборе для сравнения с $n=128$ интервал разности долей точных доставок равен $[1{,}667;11{,}667]$ процентного пункта, а разности времени на запрос --- $[6{,}467;8{,}793]$ с. Малая численность ячеек и фиксированный вычислительный контекст ограничивают интерпретацию; это не объявлено подтверждающим тестом статистического превосходства.

Неопределённость оценивалась описательным парным повторным отбором семейств внутри 20 ячеек, по три семейства в каждой, с 10\,000 повторений. Интервал разности времени на запрос между моделью и $n=256$ составил $[-5{,}286;-2{,}939]$ с. Разность долей точных доставок дала вырожденный интервал $[0;0]$, поскольку 60 пар меток совпали. Такой интервал не доказывает равенства популяционных рисков; наблюдаемые $57/60=95\%$ не являются гарантией надёжности следующего ответа $95\%$.


Даже при отсутствии наблюдаемых потерь простая консервативная верхняя оценка их безусловной доли при независимых ограниченных исходах и поправке на два сравнения равна
\begin{equation}
\overline L=\min\left\{1,\widehat L+\sqrt{\frac{\log40}{2\cdot60}}\right\}
=0{,}17533\quad\text{при }\widehat L=0.
\end{equation}
Эта оценка не доказывает неухудшение с пределом риска $5\%$. Она требует предположения независимости семейств и относится к среднему по рассматриваемому плану отбора, а не к каждому входу. Вырожденный парный интервал точности с $n=256$ также не устраняет неопределённость неизвестного популяционного риска.

\subsection*{Отдельная подвыборка Монте-Карло}
\fig{summary_05_mc20}{fig:mc20}{.96}{Время и сеточная ошибка одной реализации Монте-Карло на 20 заранее выбранных задачах. Все точки выше установленного допуска; вывод не распространяется автоматически на другие случайные реализации или параметры метода.}
Все 20 ответов конечны, но неточны. Минимальная, медианная и максимальная ошибки равны $0{,}0102624$, $0{,}0410146$, $0{,}201200$. Времена имеют минимум $9{,}612$ с, медиану $34{,}061$ с и максимум $203{,}910$ с. Суммарно Монте-Карло потребовал $1144{,}865$ с; на тех же 20 задачах Ремизов при $n=128$ --- $236{,}023$ с, при $n=256$ --- $471{,}277$ с. Стоимость на точную MC-доставку не определена при нуле таких доставок.

В одном случае Монте-Карло был быстрее $n=128$, в шести --- быстрее $n=256$, но эти ответы не достигали допуска. Поэтому среди измеренных допустимых вариантов MC-побед нет. Наблюдаемый минимум на этой подвыборке: 17 раз $n=128$ и 3 раза $n=256$. На остальных 40 уравнениях Монте-Карло не измерялся, так что полный минимум стоимости трёх допустимых вариантов для всех 60 неизвестен.

Средний прогноз вероятности успеха Монте-Карло равен примерно $0{,}06235$ при диапазоне $0{,}05261$--$0{,}09434$. Наблюдаемые $0/20$ не позволяют оценить качество распознавания положительных исходов. При одной семье на каждую из 20 ячеек внутристратификационный повторный отбор вырождается. Грубый ориентир для среднего успеха при независимых ограниченных исходах имеет вид
\begin{equation}
\left[\max(0,\widehat p-\eta),\min(1,\widehat p+\eta)\right],\quad
\widehat p=0,\quad\eta=\sqrt{\frac{\log40}{40}}\approx0{,}30368.
\end{equation}
Этот диапазон $[0;0{,}30368]$ не является интервалом надёжности нового случайного запуска на конкретном уравнении. Один повтор не разделяет дисперсию, временную дискретизацию и смещение усечения. Таким образом, подтверждена ограниченная экономия стратегии на основе Ремизова, но не способность надёжно распознавать преимущество Монте-Карло.
